\section{Methods}
\label{sec:methods}

% ~0.75 page. Code: travel-speed-gate1/t3_harness.py, t2_terrain5m.py,
% t2_terrain1m.py, t2_embed.py, t4_terrain_curve.py.
\todo{Target: the distribution of $\ln v$ for a walking trail segment.}

\todo{Covariates. (a) Slope. (b) OSM tags: sac\_scale, trail\_visibility,
surface, width, smoothness, mtb\_scale, highway, access. (c) Terrain from
elevation, read every 1~m along the segment's matched path: spread and
maximum of 5~m pitches, relief along the line of travel, side slope (5~m
grids; 1~m lidar adds STRIDE's roughness). (d) NLCD 30~m canopy and land
cover. (e) AlphaEarth satellite embeddings~\cite{alphaearth2025}.
Control: ``hidden climbing'', the extra time a fixed Tobler curve predicts
when the segment is walked as 5~m pieces.}

\todo{Models: published Tobler, Campbell (median)~\cite{campbell2022variability}
and STRIDE~\cite{campbell2024stride}; a refitted rounded asymmetric slope
curve (the baseline); additive ridge offsets on binned features; gradient
boosted trees.}

\todo{Predictive distribution: prediction plus residual quantiles from held
out training outings, so calibration has to transfer too.}

\todo{Splits (never by segment): 5 folds of outings; 5 folds of 10~km
blocks; leave one region out (17 folds). Score: CRPS~\cite{gneiting2007strictly}
as skill vs the slope curve, 95\% intervals resampling whole regions;
coverage of 90\% intervals; moving-time error per hike. References: the
hiker's pace from their first km; each region's own offset.}

\begin{equation}
\ln v = a_0 + a_1 z - (k_0 + k_1 z)\,|s + s_0|
\label{eq:terrain}
\end{equation}
with $z$ the SD of 5~m pitches along the path divided by 0.10.
\todo{Introduce Eq.~\eqref{eq:terrain}, the terrain-adjusted hiking
function: Tobler's form, with peak speed and slope sensitivity depending on
one DEM-derived number.}
